% ------------------------------------------------------------------------------
Here are useful job-management commands: 

\begin{itemize}
%\item
%\texttt{qsub ./<myscript>.sh}: once that your job script is ready,
%on \texttt{speed-submit} you can submit it using this
\item
\texttt{sbatch -A <ACCOUNT> --t <MINUTES> --mem=20G -p <PARTITION> ./<myscript>.sh}: once that your job script is ready,
on \texttt{speed-submit} you can submit it using this

\item
%\texttt{qstat -f -u <ENCSusername>}: you can check the status of your job(s)
\texttt{squeue -u <ENCSusername>}: you can check the status of your job(s)

\item
%\texttt{qstat -f -u "*"}: display cluster status for all users. 
\texttt{squeue}: display cluster status for all users. 
\option{-A} shows per account (e.g., \texttt{vidpro}, \texttt{gipsy},
\texttt{speed1}, \texttt{ai2}, \texttt{aits}, etc.),
\option{-p} per partition (\texttt{ps}, \texttt{pg}, \texttt{pt}, \texttt{pa}),
and others. \texttt{man squeue} for details.

\item
%\texttt{qstat -j [job-ID]}: display job information for [job-ID] (said job may be actually running, or waiting in the queue). 
\texttt{squeue --job [job-ID]}: display job information for [job-ID] (said job may be actually running, or waiting in the queue). 

\item
\texttt{squeue -las}: displays individual job steps (for debugging
easier to see which step failed if you used \tool{srun}).

\item
\verb+watch -n 1 "sinfo -Nel -pps,pt,pg,pa && squeue -la"+: view \tool{sinfo} information and watch the queue for your job(s).

%\item
%\texttt{qdel [job-ID]}: delete job [job-ID]. 
\item
\texttt{scancel [job-ID]}: cancel job [job-ID]. 

%\item
%\texttt{qhold [job-ID]}: hold queued job, [job-ID], from running. 
\item
\texttt{scontrol hold [job-ID]}: hold queued job, [job-ID], from running. 

%\item
%\texttt{qrls [job-ID]}: release held job [job-ID]. 
\item
\texttt{scontrol release [job-ID]}: release held job [job-ID]. 

\item
%\texttt{qacct -j [job-ID]}: get job stats. for completed job [job-ID]. \api{maxvmem} is one of the more useful stats. 
\texttt{sacct -j [job-ID]}: get job stats.
%for completed job [job-ID].
\api{maxvmem} is one of the more useful stats that you can elect to display
as a format option.

\small
\begin{verbatim}
% sacct -j 2654
JobID           JobName  Partition    Account  AllocCPUS      State ExitCode
------------ ---------- ---------- ---------- ---------- ---------- --------
2654          tcsh-test         ps     speed1          1  COMPLETED      0:0
2654.batch        batch                speed1          1  COMPLETED      0:0
2654.extern      extern                speed1          1  COMPLETED      0:0
% sacct -j 2654 -o jobid,user,account,MaxVMSize,Reason%10,TRESUsageOutMax%30
JobID             User    Account  MaxVMSize     Reason        TRESUsageOutMax
------------ --------- ---------- ---------- ---------- ----------------------
2654           serguei     speed1                  None
2654.batch                 speed1    296840K             energy=0,fs/disk=1975
2654.extern                speed1    296312K              energy=0,fs/disk=343
\end{verbatim}
\normalsize

See \texttt{man sacct} or \texttt{sacct -e} for details of the
available formatting options. You can define your preferred
default format in the \api{SACCT\_FORMAT} environment variable
in your \texttt{.cshrc} or \texttt{.bashrc} files.

\item
\texttt{seff [job-ID]}: reports on the efficiency of a job's cpu and memory utilization. 
Don't execute it on RUNNING jobs (only on completed/finished jobs), efficiency statistics may be misleading.

If you define the following directive in your batch script, your ENCS email address will receive an email with seff output when your job is finished.
\small
\begin{verbatim}
#SBATCH --mail-type=ALL        
\end{verbatim}
\normalsize

Output example:
\small
\begin{verbatim}
Job ID: XXXXX
Cluster: speed
User/Group: user1/user1
State: COMPLETED (exit code 0)
Nodes: 1
Cores per node: 4
CPU Utilized: 00:04:29
CPU Efficiency: 0.35% of 21:32:20 core-walltime
Job Wall-clock time: 05:23:05
Memory Utilized: 2.90 GB
Memory Efficiency: 2.90% of 100.00 GB
\end{verbatim}
\normalsize


\end{itemize}